\subsection{A finite lower certificate for the second Weil eigenvalue}
\label{sec:finite-weil-lower}

The following bound isolates the remaining arithmetic inequality in a
finite matrix, with an explicit error for the discarded directions.
The geometric explicit formula on $[-L,L]$ is
\begin{equation}
 q(f)=2\operatorname{Re}\bigl(F(i/2)\overline{F(-i/2)}\bigr)
       +\frac1{2\pi}\int_{\mathbb R}\Psi_L(t)|F(t)|^2\,dt,
 \label{eq:weil-symbol-lower}
\end{equation}
where $F(t)=\int_{-L}^Lf(x)e^{-itx}\,dx$ and
\[
 \Psi_L(t)=\operatorname{Re}\psi\left(\frac14+\frac{it}{2}\right)
       -\log\pi-
       \sum_{\log n<2L}\frac{2\Lambda(n)}{\sqrt n}\cos(t\log n).
\]
Here $\psi=\Gamma'/\Gamma$ and $\Lambda$ is the von Mangoldt
function.  The strict prime-power cutoff is valid because the
autocorrelation vanishes at its support endpoints.  This representation
is also the starting point of Zhu's finite reduction
\cite[\S\S2--4, pp.~7--12]{zhu2026}; the estimate below uses a direct
polynomial projection error.

Put
\[
 \mathcal A_L=\sum_{\log n<2L}\frac{2\Lambda(n)}{\sqrt n},\qquad
 T=2\pi\exp(\mathcal A_L+b+1),\qquad b>0.
\]
On $L^2([-L,L])$ write $p_\pm(x)=e^{\pm x/2}$,
$g_t(x)=e^{itx}$, and $u\otimes v:f\mapsto\langle f,v\rangle u$.
Define the compact selfadjoint operator
\[
 B_T=p_+\otimes p_-+p_-\otimes p_+
       +\frac1{2\pi}\int_{-T}^{T}
                     (\Psi_L(t)-b)g_t\otimes g_t\,dt.
\]
Let $P_N$ be the $L^2$ orthogonal projection onto polynomials of degree
at most $N$, expressed for example in the normalized Legendre basis,
and let $M_{L,T,N}=bI+P_NB_TP_N$ act on that finite-dimensional space.
For $N\ge1$, set
\[
 \begin{split}
 W_{L,T,b}&=\frac1{2\pi}\int_{-T}^{T}|\Psi_L(t)-b|\,dt,\\
 R_N&=e^{TL}\frac{(TL)^{N+1}}{(N+1)!},\\
 \varepsilon_{L,T,N}
   &=4LR_NW_{L,T,b}
        +8Le^L\frac{(L/2)^{N+1}}{(N+1)!}.
 \end{split}
\]

\begin{proposition}[A finite second-eigenvalue lower bound]
\label{prop:finite-weil-lower}
The second eigenvalue of the windowed Weil form satisfies
\begin{equation}
 a_2(L)\ge
       \min\{b,\lambda_2(M_{L,T,N})\}
           -\varepsilon_{L,T,N}.
 \label{eq:finite-weil-lower}
\end{equation}
The entries of $M_{L,T,N}$ and an upper bound for its error term use
only the displayed finite prime-power sum and integrals over finite
intervals.
\end{proposition}

\begin{proof}
We first justify the frequency cutoff, including its numerical
constant.  The elementary Euler summation formula applied to
$(x+z)^{-1}$ gives, for $\operatorname{Re}z>0$,
\[
 \psi(z)=\log z-\frac1{2z}
        +\int_0^\infty\frac{\{x\}-1/2}{(x+z)^2}\,dx.
\]
For example, apply integration by parts on each unit interval, sum up
to an integer $M$, and take the limit in
$\psi(z)=\lim_{M\to\infty}(\log M-\sum_{j=0}^{M-1}(j+z)^{-1})$.
At $z=1/4+it/2$, $t\ne0$, the integral in absolute value is at most
\[
 \frac12\int_0^\infty
     \frac{dx}{(x+1/4)^2+t^2/4}
       \le\frac{\pi}{2|t|}.
\]
Together with $|1/(2z)|\le1/|t|$ and
$\log|z|\ge\log(|t|/2)$, this proves
\[
 \Psi_L(t)\ge\log\frac{|t|}{2\pi}
                     -\frac3{|t|}-\mathcal A_L.
\]
The right side increases with $|t|$.  At our choice of $T$ it is
$b+1-3/T>b$, since $T\ge2\pi e$.  Thus
$\Psi_L(t)\ge b$ outside $[-T,T]$.

Parseval's identity in \eqref{eq:weil-symbol-lower} now gives the
quadratic-form inequality
\[
 q(f)\ge b\|f\|_2^2+\langle B_Tf,f\rangle.
\]
It first holds on the compact smooth core and then on the closed form
domain.  Notice that the two signs of the pole term are both present;
on even functions it is positive and on odd functions it is negative.

The degree-$N$ Taylor polynomial approximates $g_t$ uniformly on
$[-L,L]$ with error at most
$e^{|t|L}(|t|L)^{N+1}/(N+1)!$.  Orthogonal projection is the best
$L^2$ approximation, so
\[
 \|(I-P_N)g_t\|_2\le\sqrt{2L}R_N,\qquad
 \|g_t\|_2=\sqrt{2L}\quad (|t|\le T).
\]
Writing $r=(I-P_N)g_t$ and using
\[
 g_t\otimes g_t-P_Ng_t\otimes P_Ng_t
      =r\otimes g_t+P_Ng_t\otimes r
\]
shows that the operator norm of this difference is at most $4LR_N$.
Its integral contribution is therefore at most $4LR_NW_{L,T,b}$.

The same approximation with $t=\pm i/2$ gives
\[
 \|p_\pm\|_2\le\sqrt{2L}e^{L/2},\qquad
 \|(I-P_N)p_\pm\|_2
       \le\sqrt{2L}e^{L/2}\frac{(L/2)^{N+1}}{(N+1)!}.
\]
Expand each of the two mixed rank-one terms just as above.  Their
total error is at most the second summand of
$\varepsilon_{L,T,N}$.  We have proved
\[
 \|B_T-P_NB_TP_N\|\le\varepsilon_{L,T,N}.
\]
On the orthogonal complement of the polynomial space,
$bI+P_NB_TP_N$ equals $bI$.  Its second min--max value is consequently
$\min\{b,\lambda_2(M_{L,T,N})\}$.  The min--max principle applied to
the form inequality and this operator norm estimate proves
\eqref{eq:finite-weil-lower}.
\end{proof}

There is a stronger bound which keeps the coupling quadratic in its
size.  In the notation above put
\[
 \begin{split}
 a&=\lambda_2(M_{L,T,N}),\qquad
 r=R_N,\qquad s=\frac{(L/2)^{N+1}}{(N+1)!},\\
 c&=b-2Lr^2W_{L,T,b}-4Le^Ls^2,\\
 z&=2LrW_{L,T,b}+4Le^Ls.
 \end{split}
\]

\begin{proposition}[The coupling retained in the lower certificate]
\label{prop:finite-weil-block}
For $N\ge1$,
\begin{equation}
 a_2(L)\ge\frac{a+c-\sqrt{(a-c)^2+4z^2}}2.
 \label{eq:finite-weil-block}
\end{equation}
For a target $u<\min\{a,c\}$, the right side exceeds $u$ precisely
when $(a-u)(c-u)>z^2$.
\end{proposition}

\begin{proof}
Write $Q_N=I-P_N$.  The same bounds on $g_t$ give
\[
 \|Q_N(g_t\otimes g_t)Q_N\|\le2Lr^2,
 \qquad
 \|Q_N(g_t\otimes g_t)P_N\|\le2Lr.
\]
The two pole ranks contribute at most $4Le^Ls^2$ to the first block
and $4Le^Ls$ to the second.  Thus the tail of $bI+B_T$ is bounded
below by $cI$, and its mixed block has norm at most $z$.
Exclude one unit eigenvector belonging to the smallest eigenvalue of
$M_{L,T,N}$.  This defines a codimension-one subspace of the full
Hilbert space, and on it the head block is bounded below by $aI$.
For $f=x+y$ in that subspace, with $x=P_Nf$ and $y=Q_Nf$,
\[
 q(f)\ge a\|x\|_2^2+c\|y\|_2^2-2z\|x\|_2\|y\|_2.
\]
The minimum of this two-variable quadratic form on
$\|x\|_2^2+\|y\|_2^2=1$ is the smaller eigenvalue of
$\left(\begin{smallmatrix}a&-z\\-z&c\end{smallmatrix}\right)$.
The codimension-one min--max characterization proves
\eqref{eq:finite-weil-block}.  Positivity of this matrix after
subtracting $uI$ is exactly the stated pair of positive diagonal
conditions and determinant inequality.  No positivity of $a$ or $c$
is needed for the lower bound itself.
\end{proof}

For computations the Taylor bounds can be replaced by smaller
Legendre tails.  Define, whenever its denominator is positive,
\[
 \rho_N(v)^2=
 \frac{(2N+3)v^{2N+2}}{((2N+3)!!)^2}
 \left(1-\frac{v^2}{(2N+3)(2N+5)}\right)^{-1}.
\]
Both lower bounds remain valid with
$r=\rho_N(TL)$ and $s=\rho_N(L/2)$, using
$\varepsilon=4LrW_{L,T,b}+8Le^Ls$ in
\eqref{eq:finite-weil-lower}.  To prove this, the normalized Legendre
basis $p_n(x)=\sqrt{(2n+1)/(2L)}P_n(x/L)$ has Fourier coefficients
\[
 \langle g_t,p_n\rangle
        =\sqrt{2L(2n+1)}\,i^n j_n(tL).
\]
Rodrigues' formula, integrated by parts $n$ times, gives
\[
 j_n(v)=\frac{v^n}{2^{n+1}n!}
                  \int_{-1}^1e^{ivu}(1-u^2)^n\,du.
\]
The beta integral yields
$|j_n(v)|\le |v|^n/(2n+1)!!$ for real $v$; for imaginary $v$ the
same bound acquires at most the factor $e^{|v|}$.
The successive squared terms
$(2n+1)v^{2n}/((2n+1)!!)^2$ have ratio
$v^2/((2n+1)(2n+3))$.  Summing from $n=N+1$ bounds
$\|(I-P_N)g_t\|_2^2$ by $2L\rho_N(TL)^2$.
For $p_\pm$, the corresponding bound is
$2Le^L\rho_N(L/2)^2$.  These are exactly the two estimates used in
the block proof.

The remaining positive-power comparison can now be specified without
an unspecified approximation error.  In the parameters $L=\log\lambda$
and $\beta=9+\alpha$, its finite arithmetic quantity is
\[
 \min\{b,\lambda_2(M_{L,T,N})\}-\varepsilon_{L,T,N}.
\]
A lower bound of size $\lambda^{9+\alpha}e^{-4\pi\lambda^2}$
with one $\alpha>0$ along an unbounded sequence of windows is the
inequality still needed in this reduction.  The prolate upper estimate
then makes the corresponding Rayleigh ratios tend to zero along that
sequence.  No such matrix inequality for an unbounded sequence is
proved here.  The large value of $T$ records the cost of bounding the
prime sum by its total mass; a sharper estimate must use cancellation
or the way the translations act on functions supported in the window.
